\documentclass[12pt]{article}


\usepackage{mathtools}
\usepackage{booktabs}
\usepackage{enumitem}
\usepackage{microtype}
\usepackage{setspace}
\usepackage{hyperref}
\usepackage{natbib}
\usepackage{amsthm}
\usepackage{lastpage}

% Acknowledgments environment (standalone, not from jmlr2e)
\newcommand{\acks}[1]{\paragraph{Acknowledgments}#1}

\newcommand{\bb}{\mathbf}
\newcommand{\eps}{\varepsilon}
\newcommand{\KL}{\operatorname{KL}}
\newcommand{\E}{\mathbb{E}}
\newcommand{\R}{\mathbb{R}}
\newcommand{\Prob}{\mathbb{P}}
\newcommand{\Var}{\operatorname{Var}}
\newcommand{\Iden}{\operatorname{I}}
\newcommand{\PiW}{\Pi_{W}}
\newcommand{\argmin}{\operatorname*{arg\,min}}

% Theorem environments (amsthm)
\newtheorem{theorem}{Theorem}[section]
\newtheorem{lemma}[theorem]{Lemma}
\newtheorem{proposition}[theorem]{Proposition}
\newtheorem{corollary}[theorem]{Corollary}
\newtheorem{conjecture}[theorem]{Conjecture}
\newtheorem{definition}[theorem]{Definition}
\newtheorem{assumption}[theorem]{Assumption}
\newtheorem{prediction}[theorem]{Prediction}
\theoremstyle{remark}
\newtheorem{remark}{Remark}

% Keywords environment (standalone, not from jmlr2e)
\newenvironment{keywords}{\noindent\textbf{Keywords:}}{}




\title{A Rate--Distortion Framework for Language-Model Scaling Laws:
Entropy Floor, Two-Bottleneck Exponents, and the Joint Law}
% TODO (author block): fill in author name(s), affiliation(s), and
% corresponding-author email before submission.  No information below is
% real; every field is a placeholder requiring manual completion.
\author{%
  \name \textsc{[Author Name]} \email \textsc{[corresponding@example.org]}\\
  \addr \textsc{[Affiliation, Country]}
  \AND
  \name \textsc{[Author Name]} \email \textsc{[author@example.org]}\\
  \addr \textsc{[Affiliation, Country]}
}


\begin{document}

\begin{center}{\LARGE\bf Scaling Laws via Rate-Distortion Theory}\end{center}\doublespacing

\hypersetup{
  pdftitle={The Rate--Distortion Theory of Language-Model Scaling Laws: Entropy Floor, Two-Bottleneck Exponents, and the Joint Law},
  pdfauthor={[Author Names — fill before submission]},
  pdfsubject={},
  pdfkeywords={},
  pdfcreator={},
}

\maketitle

% ------------------------------------------------------------------
\begin{abstract}%   <- trailing '%' for backward compatibility of .sty file
The empirical scaling law
$L(N,D) = E + A N^{-\alpha_N} + B D^{-\alpha_D}$ for next-token log-loss is
the central quantitative regularity of modern language models, yet prior
theory either fits $E$ as a free parameter or derives upper bounds without
identifying the irreducible term. Using
the exact rate--distortion identity $R(D)=H(Y|X)-D$ for log-loss, we develop
such a theory. We show (i) an entropy-floor decomposition identifying $E$ as
the conditional entropy, a computable quantity; (ii) a capacity-limited
model-size exponent $\alpha_N=\gamma(2s/\beta_{\rm reg}-1)$, with a
boundary-degeneracy correction for near-deterministic contexts, an exact
closed-form toy-model mechanism whose extrapolation to real
code-versus-prose scaling (Prediction~\ref{pred:boundary}) our
pre-registered ladder test contradicted; (iii) a data-limited exponent
$\alpha_D=\gamma_{\rm ent}/(2\beta_{\rm corr})$ from temporal corpus
statistics, subsuming Cagnetta et al. [5]; and (iv) a joint law with a vanishing
cross term and self-truncating variance, yielding two compute-optimality
predictions, a near-linear $N^*\propto D^{\alpha_D/\alpha_N}$ tradeoff
(consistent with the Chinchilla law) and a superlinear
$N^*\propto D^{1/(\alpha_N+\gamma)}$ tradeoff. In principle, no exponent is a
free parameter: each is determined by corpus statistics. We report an empirical
investigation of the spectral quantities entering these formulas. Bootstrap
analysis of the token-covariance spectrum shows that the eigenvalue decay is
range-dependent rather than a clean power law, the channel-smoothness quantity
$S_i^2$ is non-monotonic, the conditional-entropy exponent $\gamma_{\rm ent}$
remains unresolvable with current estimators, and $\gamma$ (the effective-mode
scaling exponent) has not been measured for any real architecture. A
synthetic-control experiment---applying the same pipeline to data with known
power-law structure---shows the estimator fails to recover the designed
slopes, indicating the real-corpus spectral failure is partly
measurement-confounded. The data-limited prediction
$\alpha_D=\gamma_{\rm ent}/(2\beta_{\rm corr})$ therefore cannot be
non-circularly evaluated at present. The entropy-floor identification, the
data-limited recovery structure, and the joint law do not depend on
Prediction~\ref{pred:boundary} and are unaffected by its test outcome.
\end{abstract}

\begin{keywords}
scaling laws; language models; rate--distortion; cross-entropy; minimax
estimation
\end{keywords}

% ------------------------------------------------------------------
\section{Introduction}
\label{sec:intro}

The cross-entropy loss of language models trained on $D$ tokens with $N$
parameters is fit over many orders of magnitude by
\begin{equation}
  L(N,D) = E + A N^{-\alpha_N} + B D^{-\alpha_D},
  \label{eq:law}
\end{equation}
with $\alpha_N\!\approx\!0.34$, $\alpha_D\!\approx\!0.28$
[1, 2].  The exponents are fitted and $E$ is a free
constant.  We develop a first-principles theory for \eqref{eq:law} based on
the log-loss rate--distortion identity $R(D)=H(Y|X)-D$
[3, 4], from which three things follow: (i)~the floor $E$
is \emph{identified} as the conditional entropy $H(Y|X)$
(Theorem~\ref{thm:identified}); (ii)~the exponents are expressed in terms of
corpus statistics (Theorems~\ref{thm:modelN} and \eqref{eq:alphaD});
(iii)~the additive structure is derived rather than assumed
(Theorems~\ref{thm:additive}--\ref{thm:saturation}).

Figure~\ref{fig:overview} provides a schematic of the framework.

\begin{figure}[t]
\centering
\includegraphics[width=0.75\textwidth]{../results/figures/Figure_Overview_Framework.pdf}
\caption{Overview of the rate--distortion framework and empirical evidence
status. Corpus statistics enter the exponents under explicit assumptions;
empirical validation identifies partially supported, contradicted,
unresolved, and measurement-confounded components.}
\label{fig:overview}
\end{figure}

\subsection*{Contributions}
This paper makes four contributions, two theoretical and two methodological:
\begin{enumerate}[label=\textbf{(\roman*)}]
\item \textbf{Entropy-floor identification.}
  The exact rate--distortion identity for log-loss yields the decomposition
  $L(\theta)=H(Y|X)+\overline{\KL}(\theta)$, identifying the fitted floor $E$
  as the conditional entropy---a data quantity, not a free parameter
  (Theorem~\ref{thm:decomp}).
\item \textbf{Conditional two-bottleneck exponents.}
  Under explicit spectral and smoothness assumptions, the model-size exponent
  $\alpha_N=\gamma(2s/\beta_{\rm reg}-1)$ and data-size exponent
  $\alpha_D=\gamma_{\rm ent}/(2\beta_{\rm corr})$ are derived, not fitted, and
  the joint additive law follows from a vanishing cross term and variance
  saturation (Theorems~\ref{thm:modelN}--\ref{thm:saturation}).
\item \textbf{Joint law and compute-optimality consequences.}
  The two-bottleneck decomposition yields two distinct compute-optimal regimes:
  near-linear $N^*\!\propto\!D^{\alpha_D/\alpha_N}$ (source-limited) and
  superlinear $N^*\!\propto\!D^{1/(\alpha_N+\gamma)}$ (resolution-limited), with
  a crossover locus determined by the same corpus statistics
  (Section~\ref{sec:joint}).
\item \textbf{Empirical identifiability audit.}
  We systematically test the corpus quantities required by the theory
  ($\beta_{\rm reg}$, $s$, $\gamma_{\rm ent}$, $\gamma$) and report: partial
  support for the spectral structure, contradiction of the boundary-degeneracy
  extrapolation, non-identifiability of $\gamma_{\rm ent}$, and a
  synthetic-control experiment showing the spectral estimator is partly
  measurement-confounded
  (Sections~\ref{sec:empirical}--\ref{sec:prediction}).
\end{enumerate}
The empirical contribution includes identifying where the theory can and
cannot currently be instantiated.

Our main results are as follows.
\begin{enumerate}[label=\textbf{(\arabic*)}]
\item \textbf{Entropy-floor decomposition (Theorem~\ref{thm:decomp}).}
  $L(\theta) = H(Y|X) + \overline{\KL}(\theta)$,
  so $E=H(Y|X)$ is computable from the data distribution.

\item \textbf{Model-size exponent (Theorem~\ref{thm:modelN}).}
  $\alpha_N = \gamma(2s/\beta_{\rm reg}-1)$ for interior channels (Case~I),
  with a slower boundary exponent (Case~II, Proposition~\ref{prop:boundary}).
  The Case~II extrapolation to real domains is pre-registered but contradicted
  in its first test (Section~\ref{sec:prediction}).

\item \textbf{Data-limited exponent (Section~\ref{sec:data}).}
  $\alpha_D = \gamma_{\rm ent}/(2\beta_{\rm corr})$, subsuming
  Cagnetta et al. [5] as the data face of the joint law.

\item \textbf{Joint law (Theorems~\ref{thm:additive}--\ref{thm:saturation}).}
  The cross term vanishes; the variance saturates at
  $W^*(D)\!\sim\!D^{\beta_{\rm reg}/(2s)}$, yielding near-linear
  $N^*\!\propto\!D^{\alpha_D/\alpha_N}$ (source-limited) and superlinear
  $N^*\!\propto\!D^{1/(\alpha_N+\gamma)}$ (resolution-limited)
  compute-optimal predictions.
\end{enumerate}

In principle, no exponent is free: each is determined by corpus statistics. In
practice, several spectral quantities ($\gamma$, $s$, $\gamma_{\rm ent}$) have
not been reliably measured (Section~\ref{sec:discussion}).

\paragraph{What this paper does not claim.}
We do not predict prefactors $A,B$; only exponents, the floor, and their
relations.  We do not claim the trained model is rate--distortion optimal; where
the optimization gap is uncontrolled, results degrade to bounds.

The central contribution is therefore not a new fitted scaling curve, but a
conditional rate--distortion framework that derives the model- and data-limited
exponents jointly and then tests whether the corpus quantities required by those
derivations are empirically identifiable.

% ------------------------------------------------------------------
\section{Related work}
\label{sec:related}

\paragraph{Empirical scaling laws.}
Kaplan et al. [2] established the power-law fit
$L(N,D)=E+AN^{-\alpha_N}+BD^{-\alpha_D}$ and noted that $E$ corresponds to
the entropy, without deriving this identification from an information identity.
Hoffmann et al. [1] refined the exponents and the compute-optimal frontier;
both fit $E$ as a free parameter.

\paragraph{Data-limited corpus-statistics theory.}
Cagnetta et al. [5] derive the data-limited exponent
$\alpha_D=\gamma_{\rm ent}/(2\beta_{\rm corr})$ from temporal corpus statistics
and a telescoping decomposition of per-horizon losses. We recover their result
as the data face of the joint law (Theorem~\ref{thm:recovery}).

\paragraph{Spectral, kernel, and regression scaling.}
Bi and Calhoun [6] obtain the closed-form data-size regression exponent
$\alpha=2s/(2s+1/\beta)$ from a bias--variance analysis; their $\alpha$ is a
regression exponent, not the log-loss model-size exponent
$\alpha_N=\gamma(2s/\beta_{\rm reg}-1)$ derived here
(Remark~\ref{rem:registry}). Bahri et al. [7] develop a four-regime taxonomy
from kernel/NTK analyses. Bordelon et al. [8] model SGD dynamics and
derive spectral exponents. None of these frameworks identifies the entropy
floor from an information identity, derives the joint additive law, or
analyzes boundary degeneracy.

\paragraph{Information-theoretic bounds.}
Jeon and Van Roy [9] derive information-theoretic upper bounds on prediction
error; their bounds do not identify the floor or derive a data-size exponent.
Merity et al. [10] identify $E$ with the unigram entropy,
which underestimates the true floor $H(Y|X)$.

\paragraph{Gap and this paper's contribution.}
Prior frameworks fall into four categories---empirical phenomenological fits,
data-limited corpus-statistics theory, spectral/kernel/regression exponents,
and optimization-dynamics models---but none simultaneously (i)~derives the
entropy floor from an information identity, (ii)~obtains conditional
model-size and data-size exponents within one framework, (iii)~derives the
joint additive law and its compute-optimal consequences, and
(iv)~tests whether the corpus quantities entering those exponents are
empirically identifiable.  This paper addresses all four.
Figure~\ref{fig:positioning} summarizes the conceptual positioning.

\begin{figure}[t]
\centering
\includegraphics[width=0.95\textwidth]{../results/figures/Figure_Prior_Work_Positioning.pdf}
\caption{Conceptual positioning of the present rate--distortion framework
relative to major prior approaches to language-model scaling.  The present
work combines an entropy-floor identity, conditional model- and data-limited
exponents, a joint law, and an explicit empirical identifiability assessment.}
\label{fig:positioning}
\end{figure}

% ------------------------------------------------------------------
\section{Preliminaries and the entropy-floor identity}
\label{sec:identity}

\subsection{Formal setting}
Let $(X,Y)\sim\Prob$ be a next-token source: $X$ is a context, $Y$ a token
from a finite alphabet $\mathcal{Y}$, and $\Prob_{Y|x}$ the conditional law.
Write $H\equiv H(Y|X)=\E[\log 1/\Prob(Y|X)]$ for the conditional entropy. A
model is a distribution-valued map $f_\theta:x\mapsto\hat p_\theta(\cdot|x)$
with $|\theta|=N$ parameters, trained on $D$ i.i.d.\ samples. The population
loss is $L(\theta)=\E_{x,y}[-\log \hat p_\theta(y|x)]$, and the mean-KL excess
is
$\overline{\KL}(\theta)=\E_x[\KL(\Prob_{Y|x}\|\hat p_\theta(\cdot|x))]$.

\begin{assumption}[A1, log-loss]
The distortion is log-loss, $d(y,\hat p)=-\log \hat p(y)$.
\end{assumption}
\begin{assumption}[A2, finite entropy]
$H(Y|X)<\infty$, and the source admits the spectral structure of
Assumption~\ref{ass:spectral} below.
\end{assumption}

\begin{theorem}[Entropy-floor decomposition]
\label{thm:decomp}
Under A1--A2,
\begin{equation}
  L(\theta) = H(Y|X) + \overline{\KL}(\theta),
  \label{eq:decomp}
\end{equation}
for every $\theta$ in the model family.
\end{theorem}

\begin{proof}
$L(\theta)=\E[-\log\hat p_\theta]=\E[-\log\Prob]+\E[\log(\Prob/\hat
p_\theta)]=H(Y|X)+\overline{\KL}(\theta)$.
\end{proof}

\begin{corollary}
\label{cor:floor}
$L(\theta)\ge H(Y|X)$ for every $\theta$, with equality iff
$\hat p_\theta(\cdot|x)=\Prob_{Y|x}$ for $\Prob_X$-a.e.\ $x$. The loss can
therefore not be driven below the conditional entropy by any model, regardless
of $N$, $D$, or compute.
\end{corollary}

\begin{corollary}[Reframing]
\label{cor:excess}
Writing $\eps(\theta)=L(\theta)-H(Y|X)=\overline{\KL}(\theta)\ge0$, the
empirical law \eqref{eq:law} is the statement that $\eps$ decays as
$A N^{-\alpha_N}+B D^{-\alpha_D}$; it is \emph{not} a statement about
an information-theoretic capacity deficit.
\end{corollary}

\begin{theorem}[Rate--distortion tightness of the floor]
\label{thm:rd}
For a conditional source $(X,Y)$ with side information $X$ available to encoder
and decoder, log-loss distortion, and $H(Y|X)<\infty$, the rate--distortion
function is
\begin{equation}
  R(D)=H(Y|X)-D,\qquad 0\le D\le H(Y|X).
  \label{eq:rd}
\end{equation}
\end{theorem}
\begin{remark}
Theorem~\ref{thm:rd} is classical [3, 4]. It upgrades
Corollary~\ref{cor:floor} from ``our model family cannot beat $H$'' to
``no coding scheme can beat $H$.''
\end{remark}

\subsection{The floor is identified}
\begin{assumption}[A3$'$, mean-KL consistency]
\label{ass:consistent}
For the training procedure, $\overline{\KL}(\theta_{N,D})\to0$ as
$N,D\to\infty$.
\end{assumption}

\begin{theorem}[Identified floor]
\label{thm:identified}
Under A1, A2, and A3$'$, the fitted constant $E$ in \eqref{eq:law} satisfies
$E=H(Y|X)$: it is the conditional entropy, computable from the data
distribution, and is not a free parameter.
\end{theorem}

\begin{proof}
From Corollary~\ref{cor:excess}, $L-\eps=H$; by A3$'$, $\eps\to0$, so any
consistent fit returns $E=H$.
\end{proof}

\begin{remark}[Per-context-length statement]
\label{rem:perlength}
The floor is context-length dependent: $E_k=H(Y|X_{1:k})$, which decays in $k$
[11].  Finite-context models absorb the residual as inflated
$E$.
\end{remark}

\begin{remark}[Where optimization slack enters]
\label{rem:slack}
Training minimizes empirical loss, so
$\overline{\KL}(\theta_{N,D})=[L-L_{\min}(D)]+[L_{\min}(D)-H]$ splits into an
optimization gap and an estimation excess.  The theory derives the estimation
excess; the optimization gap must be controlled or absorbed
(Assumption~\ref{ass:fastlearn}).
\end{remark}

% ------------------------------------------------------------------
\section{Assumptions and notation}
\label{sec:assumptions}

\begin{assumption}[A4, spectral capacity]
\label{ass:spectral}
The token covariance has eigenvalues $\lambda_i\sim i^{-1/\beta_{\rm reg}}$
($\beta_{\rm reg}>1$), and the channel is smooth of index $s$:
$\Prob_{y}=K_X^s v_y$ with $\sup_y\|v_y\|^2<\infty$.  The model with $N$
parameters realizes the projection onto the top $W(N)\!\sim\!N^{\gamma}$
eigenmodes ($\gamma\!=\!1/2$ at fixed depth).
\end{assumption}

\begin{assumption}[A5, fast learning]
\label{ass:fastlearn}
The trained model achieves the minimax rate of the estimation problem
(up to an excess that is vanishing or absorbable as a constant), shared with
Cagnetta et al. [5].
\end{assumption}

\noindent\textbf{Notation.}
$\beta_{\rm reg}$: spectral tail index;
$s$: channel smoothness;
$\gamma$: width exponent ($W\!\sim\!N^{\gamma}$);
$\alpha_N$: model-size exponent;
$\gamma_{\rm ent}$: entropy decay ($H_n\!-\!H_\infty\!\sim\!n^{-\gamma_{\rm ent}}$);
$\beta_{\rm corr}$: token-covariance decay;
$\alpha_D$: data-size exponent.  The two ``$\beta$'' objects are distinct
(spectral vs.\ temporal).

\begin{remark}[Registry vs.\ other papers]
\label{rem:registry}
Bahri et al. [7] write $\lambda_i\sim i^{-(1+s')}$, so
$s'=1/\beta_{\rm reg}-1$. Bi--Calhoun's $\beta=\beta_{\rm reg}$;
Cagnetta et al. [5]'s $(\beta,\gamma)=(\beta_{\rm corr},\gamma_{\rm ent})$.
The Bi--Calhoun exponent $\alpha=2s/(2s+1/\beta_{\rm reg})$ is a
\emph{data-size} regression exponent, distinct from our \emph{model-size}
exponent $\alpha_N=\gamma(2s/\beta_{\rm reg}-1)$; this conflation is the most
common error in reading this literature.
\end{remark}

\begin{remark}[Capacity models: C1 vs.\ C2]
\label{rem:capacity}
The mode-counting model (A4, C1, $W\!\sim\!N^{\gamma}$) is used throughout because
it delivers the clean two-sided exponent of Theorem~\ref{thm:modelN}. The
alternative entropy-capacity model (C2, in the spirit of
Jeon and Van Roy [9]) would reappear as an upper bound on the achievable
rate; we keep C1 because it captures the resolution-limited variance that C2
cannot (an entropy budget does not resolve modes).
\end{remark}

% ------------------------------------------------------------------
\section{Model-size scaling: the capacity-limited exponent}
\label{sec:modelsize}

We study the resolution-limited excess $\eps_N = L(N,\infty)-H$ at infinite
data, where the only error is the truncation of the channel by the $W\!\approx\!
N^\gamma$ resolved modes.

\subsection{The KL/weighted-$L^2$ bridge}
\begin{lemma}[KL and weighted $L^2$]
\label{lem:bridge}
For every $x$, with $p=\Prob_{Y|x}$ and $q=\hat p(\cdot|x)$,
\begin{equation}
  \KL(p\|q) \asymp \frac12 \sum_y \frac{(q(y)-p(y))^2}{p(y)}
  \qquad (\text{small-error regime}),
  \label{eq:bridge}
\end{equation}
up to constants [12].  Log-loss measures \emph{relative} error,
weighted by $1/\Prob(y|x)$.
\end{lemma}

\begin{remark}[Conditional weighting]
\label{rem:conditional}
The weight is the \emph{conditional} $1/\Prob(y|x)$, not marginal
$1/\Prob(y)$; for binary tokens this is $1/(p(1-p))$, the inverse Bernoulli
variance.  This weight diverges at the simplex boundary (Case~II).
\end{remark}

\subsection{Interior exponent (Case I)}
\begin{theorem}[Model-size exponent, Case I]
\label{thm:modelN}
Under A1, A2, A4 with $W(N)\sim N^{\gamma}$, for channels bounded away from
the simplex boundary,
\begin{equation}
  \eps_N = L(N,\infty) - H(Y|X) \asymp A\, W^{-(2s/\beta_{\rm reg}-1)}
  = A' N^{-\alpha_N},
  \qquad \alpha_N=\gamma\left(\frac{2s}{\beta_{\rm reg}}-1\right),
  \label{eq:modelN}
\end{equation}
provided $2s/\beta_{\rm reg}>1$.
\end{theorem}

\begin{proof}
By Lemma~\ref{lem:bridge} and A4, $\eps_N\!\asymp\!\frac12\E_x\sum_y
(\hat p\!-\!p)^2/p$. The $s$-source condition and interior bound on $1/p$ give
$\eps_N\!\lesssim\!\sum_{i>W} i^{-2s/\beta_{\rm reg}}
\!\asymp\! W^{-(2s/\beta_{\rm reg}-1)}$.
Full proof in Supporting Information~S2.
\end{proof}

The exponent $\alpha_N=\gamma(2s/\beta_{\rm reg}-1)$ is the classical
width-truncation exponent of kernel approximation; the step specific to log-loss is
the \emph{conditional} weighting of Lemma~\ref{lem:bridge}, which moves the
analysis to the simplex boundary and produces Case~II.

\begin{remark}[Consistency check]
\label{rem:chinchilla}
With $\gamma=1/2$ and the measured Chinchilla $\alpha_N\approx0.34$:
$2s/\beta_{\rm reg}-1\approx0.67$, implying $s\approx1.7$ for
$1/\beta_{\rm reg}\approx2$. This is an order-of-magnitude consistency note, not a fit.
\end{remark}

\subsection{Boundary degeneracy (Case II)}
Near-deterministic contexts drive $\Prob(y|x)$ to the simplex boundary, where
the log-loss weight $1/p(y|x)$ diverges.

\begin{proposition}[Boundary mechanism]
\label{prop:boundary}
If $f(x)=\Prob(a|x)$ touches $\{0,1\}$ linearly, the boundary contribution to
the mean excess is $\eps_N^{\rm bdy}\sim\frac{c^2}{4}\log W$, so
$\eps_N^{\rm bdy}\sim W^{-2}$ in the linear-touch example with
$2s/\beta_{\rm reg}=3$---one power \emph{slower} than the interior tail.
For $V\!\ge\!2$ tokens, the boundary contribution is a sum of positive terms
with no cancellation.
For approach rate $f\sim au^p$: $p<1$ gives $\eps^{\rm bdy}\sim c^2$ (milder);
$p=1$ gives $c^2\log W$; $p>1$ gives $c^2 W^{p-1}$ (more severe).
Full proofs in Supporting Information~S2.
\end{proposition}

\subsection{Parameterization: probabilities vs.\ logits}
Real transformers parameterize logits and apply softmax.

\begin{proposition}[Logit-space amplification]
\label{prop:logit}
If the model truncates the logit field $\ell_y(x)=\log\Prob(y|x)$:
Case~I is unchanged for strictly positive channels; for zero-touching
channels the log-singularity yields $\eps_N\sim W^{-1}$, amplifying the
probability-space $W^{-2}$ slowdown.
Full proof in Supporting Information~S2.
\end{proposition}

\begin{remark}[Real channels]
\label{rem:logit}
Real channels are neither strictly positive nor exactly zero; the effective
spectral content lies between the two bounds.  Toy models: Model~B
(zero-touching, logit space) flattens toward $W^{-1}$; Model~C (positive, logit
space) gives clean Case~I ($W^{-3.08}$).
\end{remark}

% ------------------------------------------------------------------
\section{Data-limited scaling: recovery of the corpus-statistics exponent}
\label{sec:data}

The data face of the law is recovered from the data-limited theory of
Cagnetta et al. [5].  Two corpus statistics:
\begin{equation}
  H_n - H_\infty \sim n^{-\gamma_{\rm ent}},
  \qquad
  \|C(n)\|_{\rm op}\sim n^{-\beta_{\rm corr}},
  \label{eq:entropy}
\end{equation}
with the data-determined horizon $n^*(P)\sim P^{1/(2\beta_{\rm corr})}$ and
their telescoping decomposition, yields
$\alpha_D=\gamma_{\rm ent}/(2\beta_{\rm corr})$.
Published measurements (Cagnetta et al. [5]): TinyStories
$\alpha_D\!\approx\!0.19$; WikiText $\alpha_D\!\approx\!0.14$ (literature
values; we do not independently reproduce them because $\gamma_{\rm ent}$ is
not robustly identifiable).

\begin{theorem}[Recovery/identification]
\label{thm:recovery}
Within our framework the Cagnetta et al.\ [5] decomposition is identified exactly:
(i)~their $H_n$ is our per-horizon floor $H(Y|X_{1:n})$; (ii)~their excess is
the differential of our mean-KL excess; (iii)~their decomposition is the
telescoping of \eqref{eq:decomp} truncated at $n^*(P)$.  Consequently
\begin{equation}
  \alpha_D=\frac{\gamma_{\rm ent}}{2\beta_{\rm corr}},
  \label{eq:alphaD}
\end{equation}
the source-limited variance of Section~\ref{sec:joint}.
\end{theorem}

\begin{proof}[Sketch]
By Remark~\ref{rem:perlength}, $L_n=H_n+\eps_n$ for each horizon $n$.
Telescoping from $n=1$ to $n^*(P)$ collapses the per-horizon identity onto
\eqref{eq:cagnetta}.  Substituting \eqref{eq:entropy} and \eqref{eq:horizon}
into the boundary term gives \eqref{eq:alphaD}.  Full details in
Supporting Information~S2.
\end{proof}

The recovery is exact because our identity \eqref{eq:decomp} holds per
horizon and their decomposition is its telescoping.

\begin{remark}[The dominance condition]
\label{rem:dominance}
The data exponent \eqref{eq:alphaD} follows when the boundary term dominates
the estimation term.  Cagnetta et al.\ [5] state this as fast learning within the
horizon; we inherit it through Assumption~\ref{ass:fastlearn}.  Full
telescoping lemma and proof in Supporting Information~S2.
\end{remark}

\begin{remark}[Discrepancy and what it predicts]
\label{rem:discrepancy}
The directly measured $\alpha_D\!\approx\!0.14$--$0.19$ differs from
$\beta_{\rm fit}\!\approx\!0.28$ because the fitted value blends regimes
(Section~\ref{sec:joint}); the pure exponent is \eqref{eq:alphaD}.
\end{remark}

% ------------------------------------------------------------------
\section{The joint $N$--$D$ law and compute optimality}
\label{sec:joint}

\begin{theorem}[Two-source decomposition]
\label{thm:additive}
Under A1, A2, A4, the mean-KL excess admits the decomposition
\begin{equation}
  \eps(N,D)=\eps_{\rm bias}(N)+\eps_{\rm var}(N,D)+r,
  \label{eq:decom}
\end{equation}
where $\eps_{\rm bias}=\frac12\E_x\sum_y
\frac{(\PiW f-f)_y^2}{f_y}$ (a truncation error, $N$-only),
$\eps_{\rm var}=\frac12\E_x\sum_y
\frac{(g-\PiW f)_y^2}{f_y}$ (an estimation error), and, for channels bounded
away from the simplex boundary,
\begin{equation}
  \E_{\rm data}\left[\sum_y\frac{(\PiW f-f)_y\,(g-\PiW f)_y}{f_y}\right]=0,
  \qquad |r|\le C(\delta)\,\eps^{3/2},
  \label{eq:cross}
\end{equation}
where $C(\delta)$ depends only on the interior bound and
$\eps=\eps_{\rm bias}+\eps_{\rm var}$. Hence the cross term vanishes in
expectation \emph{exactly} at quadratic order, and the residual is
$o(\eps)$.
\end{theorem}

\begin{proof}
Decompose $g-f=\delta_b+\delta_v$; $\E[\delta_v]=0$ gives the cross-term
equality. Full details in Supporting Information~S2.
\end{proof}

For squared loss the same statement holds \emph{pointwise} without any
expectation: $\delta_b\perp\delta_v$ in $L^2(\Prob_X)$ by orthogonality of
Fourier modes, so $\E_x[(g-f)^2]=\E_x[\delta_b^2]+\E_x[\delta_v^2]$ exactly
(verified to float64 roundoff, $10^{-9}$, Section~\ref{sec:numerics}).

\begin{theorem}[Variance saturation / self-truncation]
\label{thm:saturation}
Consider the Gaussian sequence model of the projected channel:
observations $\hat c_i=c_i+\sigma_i z_i/\sqrt{D}$ with $|c_i|\sim
i^{-s/\beta_{\rm reg}}$, and model class $M_W$ of size $W$. The minimax MSE
risk over $M_W$ is
\begin{equation}
  \inf_{q\in M_W}\ \sup_{c}\ \E\|q-f\|^2
  = (1+o(1))\sum_{i\le W}
  \frac{c_i^2\,(\sigma_i^2/D)}{c_i^2+\sigma_i^2/D},
  \label{eq:minimax}
\end{equation}
achieved by the Wiener filter with weight $w_i=
\frac{c_i^2}{c_i^2+\sigma_i^2/D}$ on mode $i$. Consequently, for $W$ larger
than the threshold
\begin{equation}
  W^*(D)\sim \left(\frac{\sigma^2}{D}\right)^{-\beta_{\rm reg}/(2s)}
  \sim D^{\beta_{\rm reg}/(2s)},
  \label{eq:Wstar}
\end{equation}
the risk is dominated by modes $i\le W^*$ and is $W$-independent:
\begin{equation}
  \eps_{\rm var}(D)\sim \frac{W^*(D)}{D}\sim
  D^{-(1-\beta_{\rm reg}/(2s))}.
  \label{eq:varstar}
\end{equation}
The same statement transfers to the KL excess in the interior (the KL is the
weighted $L^2$ error of Lemma~\ref{lem:bridge} with a mode-independent
weight, so the exponent is unchanged).
\end{theorem}

\begin{proof}
Classical minimax risk of the Gaussian sequence model
[12, 13]; the Wiener filter achieves it, and the
threshold \eqref{eq:Wstar} is SNR$=1$.
\end{proof}

\begin{remark}[What this proves, and the caveat]
\label{rem:realizability}
The minimax floor is $W$-independent: a self-truncating model behaves as if its
resolution were $\min(W(N),W^*(D))$, making the data term $N$-independent.
ERM does \emph{not} self-truncate; Assumption~\ref{ass:fastlearn} bridges the gap.
\end{remark}

\subsection{Two regimes and the origin of the additive law}
Combining Theorem~\ref{thm:modelN}, \eqref{eq:alphaD}, and
Theorem~\ref{thm:saturation}:

\begin{equation}
  \eps(N,D) = A N^{-\alpha_N}+
  \begin{cases}
    C N^{\gamma}/D, & W(N)<W^*(D)\quad \text{(resolution-limited)}, \\
    B D^{-\alpha_D}, & W(N)>W^*(D)\quad \text{(source-limited)},
  \end{cases}
  \label{eq:regimes}
\end{equation}
with $\alpha_N=\gamma(2s/\beta_{\rm reg}-1)$ and
$\alpha_D=\gamma_{\rm ent}/(2\beta_{\rm corr})$. The additive law
\eqref{eq:law} is the source-limited regime. In the resolution-limited regime
the data term carries an explicit $N^\gamma$ prefactor, so fitting
\eqref{eq:law} there is mis-specified---a predicted finite-scale artifact of
the fitted exponents.

\begin{remark}[Empirical constraint]
\label{rem:empirical}
The fitted $\beta_{\rm fit}\approx0.28$ exceeds the pure source-limited value
$\alpha_D\approx0.14$--$0.19$; regime blending explains the discrepancy.
\end{remark}

\subsection{Variance at fixed resolution and the iid source-limited exponent}
With iid contexts, the empirical Fourier coefficients $\hat c_i$ have
$\Var(\hat c_i)\approx\sigma_i^2/D$ with $\sigma_i^2\approx\E[f(1-f)]$
uniform in $i$ (verified to $<3\%$).  The resolution-limited variance is
therefore $\eps_{\rm var}\sim W/D$ (i.e., $C N^{\gamma}/D$). Self-truncation
(Theorem~\ref{thm:saturation}) freezes the effective resolution at
$W^*(D)\sim D^{\beta_{\rm reg}/(2s)}$, giving the iid source-limited exponent
$\alpha_D^{\rm iid}=1-\beta_{\rm reg}/(2s)$. For the toy numbers this is
$\approx0.85$; reaching the empirical $\beta_{\rm fit}\approx0.28$ requires
heavier tails, rougher channels, or the correlated-source mechanism
\eqref{eq:alphaD}.

\subsection{Compute-optimality: two distinct predictions}
Minimizing $L(N,D)$ under the compute budget $C=6ND$ (Chinchilla) by
Lagrange multipliers gives two crisp, distinct regimes (derivation in
Appendix~\ref{app:compute}):

\paragraph{Source-limited (additive) regime.}
$\eps=A N^{-\alpha_N}+B D^{-\alpha_D}$, $\text{ND}=C$:
\begin{equation}
  N^* \propto C^{\alpha_D/(\alpha_N+\alpha_D)},\qquad
  D^* \propto C^{\alpha_N/(\alpha_N+\alpha_D)},
  \qquad\text{i.e.}\qquad
  N^* \propto D^{\alpha_D/\alpha_N}.
  \label{eq:sourceopt}
\end{equation}
With measured $\alpha_N\approx0.34,\ \alpha_D\approx0.28$, the slope is
$\alpha_D/\alpha_N\approx0.82\approx1$: the celebrated \emph{near-linear}
tradeoff, \emph{derived} here as a consequence of $\alpha_D\approx\alpha_N$
rather than assumed. Quantitative slopes depend on exponent estimates
that remain uncertain at the current corpus/model scale.

\paragraph{Resolution-limited regime.}
$\eps=A N^{-\alpha_N}+C N^{\gamma}/D$, $\text{ND}=C$:
\begin{equation}
  N^* \propto C^{1/(\alpha_N+\gamma+1)},\qquad
  D^* \propto C^{(\alpha_N+\gamma)/(\alpha_N+\gamma+1)},
  \qquad\text{i.e.}\qquad
  N^* \propto D^{1/(\alpha_N+\gamma)}.
  \label{eq:resolopt}
\end{equation}
With $\gamma=\frac12,\ \alpha_N\approx0.34$, the slope is
$1/(\alpha_N+\gamma)\approx1.19$: \emph{superlinear}, and the two
predictions are cleanly separated (0.82 vs.\ 1.19). Their quantitative
slopes depend on $\gamma$ (not yet measured for any real architecture) and
$\alpha_N$ (estimated with uncertainty).

\begin{prediction}[Compute-optimality discriminator]
\label{pred:compute}
A measured $N^*$--$D^*$ slope $\approx1$ supports source-limited variance and
$\alpha_D/\alpha_N\approx1$; a slope strictly above $1$ supports
resolution-limited variance. The near-linear Chinchilla result is consistent
with the source-limited prediction, though the numerical slopes depend on
exponent estimates that carry substantial uncertainty.
\end{prediction}

\subsection{Crossover}
\begin{conjecture}[Crossover locus]
\label{conj:crossover}
The transition between the two regimes in \eqref{eq:regimes} occurs at
$W(N^*)\sim W^*(D)$, i.e., $N^*\sim D^{\beta_{\rm reg}/(2s\gamma)}$
($=D^{\beta_{\rm reg}/s}$ for $\gamma=\frac12$). Falsifiable signatures: the
fitted data exponent drifts toward $1$ (resolution-limited) when the model is
under-resolved relative to the data, and toward \eqref{eq:alphaD}
(source-limited) when over-resolved.
\end{conjecture}

With toy-consistency estimates ($s\!\approx\!1.7$,
$\beta_{\rm reg}\!\approx\!0.5$, $\gamma\!=\!\tfrac12$), the self-truncation
threshold is $W^*(D)\!\approx\!D^{0.15}$ and the crossover locus is
$N^*\!\approx\!D^{0.3}$: at $D=10^9$ the crossover sits near $N^*\!\approx\!500$,
well below deployed scales.  The Chinchilla slope $\approx\!1$ independently
confirms this---the source-limited prediction \eqref{eq:sourceopt} matches the
data, so the fitted frontier lives in the additive regime.  Measured
$(s,\beta_{\rm reg})$ are needed for quantitative use.

% ------------------------------------------------------------------
\section{Measuring the exponents from a corpus}
\label{sec:measurement}

The claim that no exponent is free rests on the exponents being
corpus-measurable. The spectral tail index $\beta_{\rm reg}$ is fit from the
eigenvalue decay of the token-covariance; the channel smoothness $s$ is checked
from the decay of projected next-token channels or inverted from a measured
$\alpha_N$ via Theorem~\ref{thm:modelN}. The temporal statistics $\gamma_{\rm ent}$
and $\beta_{\rm corr}$ are fit from the decay of per-horizon losses and the
operator norm of the token--token covariance at separation $n$
[5]; the consistency check is the scaling collapse of
$n$-gram loss curves. In practice, $\gamma_{\rm ent}$ is not robustly
identifiable with the current plug-in estimator (see below).

\subsection{Fitting $\alpha_N,\ \alpha_D$ and testing the floor}
A joint-law fit yields five falsifiable tests:
\begin{enumerate}[label=\textbf{(\arabic*)}]
\item \emph{Floor:} $E'$ must match $H(Y|X)$ (Theorem~\ref{thm:identified}).
\item \emph{Exponents:} fitted $\alpha_N=\gamma(2s/\beta_{\rm reg}-1)$ with
  measured $(s,\beta_{\rm reg})$; fitted data exponent drifts toward
  $\gamma_{\rm ent}/(2\beta_{\rm corr})$ as $N\!\to\!\infty$
  (Remark~\ref{rem:discrepancy}).
\item \emph{Compute slope:} $N^*$--$D^*$ slope $\approx0.82$
  (source-limited) or $\approx1.19$ (resolution-limited)
  (Prediction~\ref{pred:compute}).
\item \emph{Crossover:} data exponent drifts toward \eqref{eq:alphaD} as
  model is over-resolved (Conjecture~\ref{conj:crossover}).
\item \emph{Domain:} code shows slower $\alpha_N$ than prose
  (Prediction~\ref{pred:boundary}).
\end{enumerate}
Each test is quantitative; the theory is falsified if any fails by a
margin larger than the estimation error of the corpus statistics.

\subsection{Empirical status on \textsc{wikitext}-103}
\label{sec:empirical}

We report a first-pass measurement on \textsc{wikitext}-103-raw
[10], tokenized with BPE-1024 ($P=2.07\times10^8$ training
tokens). Conditional entropies use the plug-in estimator with Miller--Madow
correction; $\beta_{\rm corr}$ is fit from $\|C(n)\|_{\rm op}$ over lags
$9$--$3000$; $\beta_{\rm reg}$ from the eigen-spectrum of a context-covariance
window of length $4$ ($1.2\times10^6$ samples). Bootstrap CIs use $10$
subsamples of $10^6$ windows each; full results in
Appendix~\ref{app:bootstrap}.

Measured corpus statistics (BPE-1024):
$\beta_{\rm corr}\!=\!0.42$ (token), $0.40$ (char);
$\beta_{\rm reg}\!=\!5.18$ (char).  Full table in Supporting Information~S2.
The token-level decay $H_5-H_1\!\approx\!4.6$ bits over four context tokens is
too steep for the power-law model \eqref{eq:entropy} to be resolved by the
plug-in estimator; we do not report a plug-in $\gamma_{\rm ent}$.

We train a small transformer ($d=128$, two layers, $0.56$M parameters) on $D$
tokens with a fixed $12{,}000$-step budget. The loss falls from $5.01$~bits at
$D=125$K to $3.76$~bits at $D=64$M, saturating beyond $D\!\approx\!4$M. The
effective exponent depends on the fit: local slope is
$\alpha_D\!\approx\!0.10$, global floor-free slope is $\approx\!0.05$, and
fixing $E$ at plug-in floors gives $\alpha_D\!=\!0.08$--$0.29$ with bootstrap
intervals $[0.04,0.12]$--$[0.18,0.38]$ (2{,}000 resamples). The free fit gives
$E\!\approx\!3.70$, $\alpha_D\!\approx\!0.55$ (poorly conditioned
[14]). All estimates are consistent with $\alpha_D>0$ and, in
order of magnitude, with $\alpha_D\!\approx\!0.14$ from Cagnetta et al. [5].
With $\beta_{\rm corr}\!\approx\!0.42$, the fixed-$E$ profile implies
$\gamma_{\rm ent}\!\in\![0.04,0.24]$; $\gamma_{\rm ent}$ is not resolvable
from plug-in counts, and the formula $\alpha_D\!=\!\gamma_{\rm ent}/(2\beta_{\rm corr})$ cannot be non-circularly evaluated.

\subsection{Large-scale breakdown: the entropy floor is the predicted
\emph{non-scaling} component}
\label{sec:yan}

Yan et al.~[15] decompose $32$ pretrained models into three
components and report that only error-entropy follows a clean power law at
large scale; the other two components are $N$-invariant, so the log-log slope
decays. This is a \emph{prediction} of the present framework: the local slope
\begin{equation}
\frac{d\log L}{d\log N}
  = -\alpha_N\,\frac{A N^{-\alpha_N}}{H(Y|X)+A N^{-\alpha_N}}
  \longrightarrow 0 \quad (N\to\infty),
  \label{eq:slope}
\end{equation}
goes to zero. Two caveats: their decomposition is not the same object as
\eqref{eq:decomp} (we claim agreement on the \emph{macro-finding} only), and
the achievability gap of Assumption~\ref{ass:fastlearn} could grow at large
$N$---separating it from the entropy-floor mechanism requires a direct gap
estimate at $N\gtrsim10^9$.

% ------------------------------------------------------------------
\section{Boundary-degeneracy prediction}
\label{sec:prediction}

\begin{prediction}[Domain-dependent exponents]
\label{pred:boundary}
Domains with a large measure of near-deterministic next-token contexts
(code, arithmetic, punctuation-heavy text, chess) show a measurably slower
log-loss exponent than the Case~I value $\alpha_N=\gamma(2s/\beta_{\rm
reg}-1)$ at the same $(s,\beta_{\rm reg})$: the conditional law touches the
simplex boundary there, and log-loss decays one power slower than squared
loss (Proposition~\ref{prop:boundary}), amplified under logit
parameterization (Proposition~\ref{prop:logit}). Natural text (mostly
interior contexts) follows the Case~I exponent up to a small boundary
correction. The size of the shift scales with the measure of
near-degenerate contexts, a domain property.
\end{prediction}

\paragraph{Direct test against a real model ladder (pre-registered).}
\label{par:prediction-test}
We tested the comparative claim on a Pythia $70$M--$2.8$B ladder
(step-143000; $6.9$B does not fit the $12$~GB GPU) with fixed-E fits
$L(N)=E+AN^{-\alpha_N}$ using independent n-gram floors per corpus, and a
pre-registered verdict rule: SUPPORT iff boundary(code)$>$boundary(prose) and
$\alpha_N$(code)$<\alpha_N$(prose) with gap CI excluding zero. Corpora: $4$M-token
eval slices of wikitext-103-raw (prose) and github-code (code). Reference-model
boundary fractions: $0.096$ (prose), $0.488$ (code).

Result. The premise holds: the reference-model ($410$M) boundary fraction
(top-1 $p>0.95$) is $0.096$ (prose) vs.\ $0.488$ (code). With the
$H_4^{\rm MM}$ floor: $\alpha_N$(prose)$=0.17$,
$\alpha_N$(code)$=0.38$, gap$=-0.21$; with $H_5^{\rm MM}$:
$0.15$ vs.\ $0.27$, gap$=-0.12$. \textbf{Verdict: CONTRADICT}---code shows
the \emph{larger} $\alpha_N$ with the gap CI excluding zero in the wrong
direction. Figure~\ref{fig:alphaN_pred_obs} visualizes this.

\begin{figure}[t]
\centering
\includegraphics[width=0.92\textwidth]{../results/figures/Figure_AlphaN_Predicted_vs_Observed.pdf}
\caption{Predicted versus observed model-size scaling exponents $\alpha_N$.
Prose is partially consistent with the conditional theoretical prediction
over the relevant fitting range, whereas the code-domain result shows the
opposite ordering from Prediction~\ref{pred:boundary}. The comparison does
not constitute validation of the underlying assumptions.}
\label{fig:alphaN_pred_obs}
\end{figure}

We report this plainly. The toy-level mechanism is exact; what fails is its
extrapolation to real code-versus-prose. Post-hoc explanations include floor
proximity (the fixed floor sits close to code losses), the measured premise
(top-1 LM confidence) not tracking the theory's boundary object (simplex
distance), and Pythia models being trained on mixed web text. Status:
\emph{unresolved as stated}.

% ------------------------------------------------------------------
\section{Numerical verification}
\label{sec:numerics}

All toy models are exactly solvable: context $x\sim U[0,1]$, channels built
from the triangle wave, closed-form Fourier coefficients, orthogonal
projection onto the top $W$ modes, exact KL integration. Detailed tables for
the two-token Case~I/II verification, $V$-token generalization,
parameterization (logit-space), additivity cross terms, and saturation
diagnostics are given in Supporting Information~S2.
The key results are:
\begin{enumerate}[label=\textbf{(\roman*)}]
\item Case~I (bounded channel) preserves the MSE exponent; Case~II
  (boundary touch) shifts the KL exponent by $+1$ in $W$, verifying
  Proposition~\ref{prop:boundary}.
\item The Case~II shift survives $V=3$ with $\rho$-independent exponent
  $-2.00$, confirming no cancellation among positive weight terms.
\item Logit-space parameterization brackets the real situation between
  models~B and~C (Remark~\ref{rem:logit}).
\item The squared-loss cross term vanishes to $O(10^{-9})$ in exact
  arithmetic, and the KL/MSE ratio is $\approx 8$ (constant in $W$),
  consistent with the weighted-$L^2$ bridge (Lemma~\ref{lem:bridge}).
\item The self-truncating (Wiener) estimator saturates in $W$ for both
  bounded and degenerate channels; unregularized ERM does not.
\end{enumerate}

\begin{figure}[h]
\centering
\includegraphics[width=0.72\textwidth]{../results/figures/fig_synth_control.pdf}
\caption{Synthetic-control spectral diagnostics. The eigendecomposition pipeline
fails to recover the designed power-law slopes.}
\label{fig:synth}
\end{figure}

\paragraph{Synthetic-control spectral test.}
\label{numerics:synth}
We apply the identical pipeline to synthetic data with known power-law
structure (Regime~A: slope $-4.0$; Regime~B: slope $-1.0$; 3 sizes $\times$ 2
seeds). The estimator fails: Regime~A measures $-0.92$, Regime~B overshoots at
$-1.55$, and $S_i^2$ is non-monotonic. Real-corpus spectral findings are
therefore partly measurement-confounded.

% ------------------------------------------------------------------
\section{Discussion}
\label{sec:discussion}

Both exponents in \eqref{eq:law} are corpus-measurable in principle:
$\alpha_N=\gamma(2s/\beta_{\rm reg}-1)$ and
$\alpha_D=\gamma_{\rm ent}/(2\beta_{\rm corr})$; the floor is
$H(Y|X)$ (Theorem~\ref{thm:identified}).  In practice, $\gamma$ is unmeasured,
$\gamma_{\rm ent}$ is unresolvable, and the synthetic-control experiment (Figure~\ref{fig:synth}) shows
the eigenvalue-decay slope is partly measurement-confounded.  The boundary test
returned CONTRADICT; the toy mechanism is exact but its extrapolation is
unresolved.  The published fits ($\alpha_N\!\approx\!0.34$,
$\beta_{\rm fit}\!\approx\!0.28$) are consistent with the theory: $\alpha_N$ is
the Case~I exponent, and $\beta_{\rm fit}$ lies between the source-limited
$\alpha_D\!\approx\!0.14$--$0.19$ and $1$, as regime blending predicts.

Assumption~\ref{ass:fastlearn} is verified in the toy model and shared with
Cagnetta et al. [5]; if it fails, claims degrade to bounds.  Limitations include
architecture-dependent fast-learning realizability, no clean relation between
$\beta_{\rm corr}$ and $\beta_{\rm reg}$, log factors and constants, and
spectral estimator bias.  Details in Supporting Information~S2.

% ------------------------------------------------------------------
\paragraph{Summary of theoretical and empirical evidence status.}
Because the framework contains both exact theoretical statements and
empirically testable assumptions, Figure~\ref{fig:evidence_status} separates
mathematical status from empirical status to avoid conflating derivation with
validation.

\begin{figure}[h]
\centering
\includegraphics[width=0.95\textwidth]{../results/figures/Figure_Framework_Evidence_Status.pdf}
\caption{Rate--distortion framework and empirical evidence status.
The theoretical results are conditional on explicit assumptions; the
empirical investigation identifies components that are supported, contradicted,
unresolved, or measurement-confounded.}
\label{fig:evidence_status}
\end{figure}

% ------------------------------------------------------------------
\acks{We thank the anonymous reviewers of an earlier draft for helpful comments,
and acknowledge open-source tooling (\LaTeX, PyTorch, HuggingFace datasets,
tokenizers) used for the numerical verification and empirical measurements.}

% ------------------------------------------------------------------
\paragraph{Funding}
% TODO: Replace with actual funding sources
This research did not receive any specific grant from funding agencies in the
public, commercial, or not-for-profit sectors.

% ------------------------------------------------------------------
\paragraph{Competing interests}
The authors declare that they have no competing interests.

% ------------------------------------------------------------------
\paragraph{Author contributions}
% TODO: update after author block is filled.  The roles below reflect the
% actual work recorded in the repository (git history, scripts, results).
% One contributor performed all of the following; adjust if co-authors are
% added.
\textbf{[Author Name]}: conceptualization, methodology, formal analysis,
investigation, software, data curation, writing --- original draft, writing
--- review \& editing, visualization.

% ------------------------------------------------------------------
\noindent\textbf{Data Availability Statement}:
\textsc{wikitext}-103-raw and \texttt{codeparrot/github-code} are publicly
available on HuggingFace. Parquet files are included in the repository's
\texttt{data/} directory; derived data and figure scripts are in
\texttt{results/}
(\url{https://github.com/Shahul9570/Research.git}).

\noindent\textbf{Code Availability Statement}:
Analysis scripts, experiment code, and pre-registration protocol are at
\url{https://github.com/Shahul9570/Research.git}
(Python~3.12.3, PyTorch~2.6.0+cu124).

\noindent\textbf{Pre-registration}:
The Prediction~\ref{pred:boundary} test was pre-registered
(\texttt{experiments/pythia/PRE-REGISTRATION.md}); designated SI~S1.

% ------------------------------------------------------------------
\paragraph{Supporting Information.}
Full proofs, detailed numerical tables, bootstrap CI tables, and numerical-setup description are in Supporting Information~S2.

% ------------------------------------------------------------------
\bibliographystyle{plainnat}


\begin{thebibliography}{15}

\bibitem[1]{hoffmann2022}
J.~Hoffmann, S.~Borgeaud, A.~Mensch, E.~Buchatskaya, T.~Cai, E.~Rutherford,
D.~de~Las~Casas, L.~A.~Hendricks, J.~Welbl, A.~Clark, T.~Hennigan,
E.~Noland, K.~Millican, G.~van~den~Driessche, B.~Damoc, A.~Guy,
S.~Osindero, K.~Simonyan, E.~Elsen, J.~W.~Rae, O.~Vinyals, and
L.~Sifre.
Training compute-optimal large language models.
In \textit{NeurIPS}, 2022.

\bibitem[2]{kaplan2020}
J.~Kaplan, S.~McCandlish, T.~Henighan, T.~B.~Brown, B.~Chess,
R.~Child, S.~Gray, J.~Radford, J.~Wu, and D.~Amodei.
Scaling laws for neural language models.
Preprint, 2020.

\bibitem[3]{berger1971}
T.~Berger.
\textit{Rate Distortion Theory: A Mathematical Basis for Data Compression}.
Prentice-Hall, 1971.

\bibitem[4]{cover2006}
T.~M.~Cover and J.~A.~Thomas.
\textit{Elements of Information Theory}, 2nd ed.
Wiley, 2006.

\bibitem[5]{cagnetta2026}
F.~Cagnetta, A.~Ravent\'os, S.~Ganguli, and M.~Wyart.
Deriving neural scaling laws from the statistics of natural language.
In \textit{ICML}, 2026. arXiv:2602.07488v3 (2026-07-03), verified 2026-08.

\bibitem[6]{bicalhoun2025}
T.~Bi and A.~J.~Calhoun.
Scaling laws are redundancy laws.
Preprint, 2025.

\bibitem[7]{bahri2024}
Y.~Bahri, E.~Dohmatob, J.~Schmidt, M.~Beltagy, and Y.~Zhe.
Explaining neural scaling laws.
\textit{PNAS}, 121(27):e2311878121, 2024.

\bibitem[8]{bordelon2024}
B.~Bordelon, J.~D.~Canatar, and C.~Pehlevan.
A dynamical model of neural scaling laws.
In \textit{ICML}, 2024.

\bibitem[9]{jeonvanroy2024}
H.~J.~Jeon and B.~Van~Roy.
Information-theoretic foundations for neural scaling laws.
Preprint, 2024.

\bibitem[10]{merity2017}
S.~Merity, C.~Xiong, J.~Bradbury, and R.~Socher.
Pointer sentinel mixture models.
In \textit{ICLR}, 2017.

\bibitem[11]{scheibner2025}
C.~Scheibner, L.~M.~Smith, and W.~Bialek.
Large language models and the entropy of English.
arXiv:2512.24969, 2025.

\bibitem[12]{tsybakov2009}
A.~B.~Tsybakov.
\textit{Introduction to Nonparametric Estimation}.
Springer, 2009.

\bibitem[13]{efroimovich1985}
S.~Y.~Efroimovich.
Nonparametric curve estimation from a sample.
\textit{Ann. Statist.}, 13:66--99, 1985.

\bibitem[14]{besiroglu2024}
T.~Besiroglu, E.~Erdil, M.~Barnett, and J.~You.
Chinchilla scaling: a replication attempt.
arXiv:2404.10102, 2024.

\bibitem[15]{yan2026}
J.~Yan, Z.~Wei, Q.~Ai, Y.~Liu, and J.~Zhan.
What scales in cross-entropy scaling law?
In \textit{ICLR}, 2026. arXiv:2510.04067.

\end{thebibliography}

\end{document}
